\section{Contexto e motivação}\label{sec:1.1}

A atualidade é marcada pela revolução digital, que amplia continuamente a presença de elementos tecnológicos no cotidiano das pessoas e no centro das disputas geopolíticas. Nesse cenário, o Brasil detém a segunda maior reserva de terras raras do mundo e a Agência Nacional de Mineração estima que o país detenha cerca de 14\% do total global conhecido, ainda que haja incerteza quanto ao tamanho real desse estoque~\cite{senado2026terras}. Essa posição confere ao país papel estratégico no fornecimento de minerais críticos para a fabricação de equipamentos e tecnologias em diversas áreas, tais como comunicação, energia renovável, saúde e indústria militar~\cite{senado2026terras}. A título de exemplificação, motores de carros elétricos, componentes eletrônicos, chips avançados, aceleradores de inteligência artificial e equipamentos de computação em nuvem estão entre os dispositivos que dependem desses minerais para sua fabricação.

A despeito dessa posição, a disponibilidade do minério, por si só, não se converte em desenvolvimento tecnológico, haja vista que o Brasil ainda exporta a matéria-prima bruta a baixo valor e importa o produto final de alto custo, em razão da ausência de capacidade interna de refino e beneficiamento~\cite{senado2026terras}. Nesse sentido, a necessidade não só de extrair e exportar esse insumo, mas também de processá-lo e de agregar valor é premente para a soberania do país, com o intuito de que ele ocupe o papel de competidor industrial e se desenvolva frente às tecnologias emergentes. Com efeito, é justamente na etapa de maior valor agregado da cadeia, isto é, no projeto, e não na extração, que a dependência tecnológica se manifesta com maior nitidez.

Nessa conjuntura, o domínio do projeto de circuitos integrados é uma das chaves para essa autonomia tecnológica. A fim de entender por que esse domínio importa, é preciso antes definir o conceito: o circuito integrado exerce a função de núcleo de comando dos sistemas eletrônicos, o que o torna responsável por aspectos como desempenho, confiabilidade e velocidade desses equipamentos. Também conhecido como chip ou microchip, o circuito integrado reúne, nos dispositivos atuais, bilhões de transistores em um único componente~\cite{patterson2020computer}. Em síntese, a capacidade de fabricar chips significa também a independência geopolítica de produzir os próprios processadores e microcontroladores, os quais são a essência da computação.

Por conseguinte, o projeto de circuitos integrados exige também o domínio do subsistema de memória. A memória é o componente responsável por armazenar os dados e os programas a serem executados pela máquina e está ligada diretamente ao desempenho do processador, na medida em que define a latência de acesso e a largura de banda disponíveis ao núcleo. Esse descompasso se traduz no \emph{memory wall} (gargalo de memória)~\cite{wulf1994memorywall}, fenômeno em que a CPU\glsadd{CPU} processa dados muito mais rapidamente do que a memória consegue entregá-los e, como consequência, permanece ociosa à espera de respostas. Diante desse contexto, surge o conceito de memória cache, objeto central deste trabalho: ela armazena cópias de instruções e dados que a CPU tende a reutilizar, reduzindo a necessidade de acesso à memória principal, mais lenta. A eficácia dessa estratégia depende do princípio da localidade, segundo o qual os programas tendem a acessar, em um intervalo curto de tempo, um conjunto restrito de endereços~\cite{patterson2020computer}.

\section{O projeto e o recorte individual}\label{sec:1.2}

Este trabalho integra um esforço coletivo de desenvolvimento cujo objetivo é a construção de um processador RISC-V capaz de executar o sistema operacional Linux, descrito integralmente na linguagem Chisel. A escolha da linguagem não é incidental: por ser embarcada em Scala, o Chisel permite parametrizar e reutilizar módulos de hardware com maior concisão do que as linguagens de descrição tradicionais, e o ambiente que o acompanha viabiliza a verificação automatizada com ScalaTest e ChiselTest sobre o simulador Treadle, que acompanha a biblioteca e executa o circuito elaborado dentro da própria máquina virtual, com a opção de execução sobre o Verilator em casos específicos.

Nessa perspectiva, o projeto está organizado em frentes complementares, distribuídas entre os integrantes do grupo. A primeira responde pelo núcleo de processamento, compreendendo a organização do pipeline, a unidade lógica e aritmética, o banco de registradores, o decodificador e a unidade de controle, bem como o tratamento de conflitos de dados e de controle. A segunda ocupa-se da infraestrutura de sistema necessária à inicialização do Linux, que abrange os registradores de controle e status, o mecanismo de traps, a temporização e tradução de endereços sob o esquema Sv32. Por fim, a terceira frente, objeto desta monografia, responde pela hierarquia de memória.

Em vista disso, convém delimitar com precisão o que constitui contribuição desta monografia. O pipeline de três estágios que serve de base à integração foi desenvolvido por Victor Henrique da Silva Costa, integrante do grupo, e é aqui empregado como componente consumidor, não como objeto de investigação: os estágios de busca, decodificação e execução, o banco de registradores, a unidade lógica e aritmética e o decodificador de instruções não são descritos nem avaliados neste trabalho, e sua corretude é pressuposta a partir da verificação realizada por seu autor. Assim, a contribuição aqui apresentada concentra-se na especificação, implementação e verificação do subsistema de cache, compreendendo as memórias L1 de instrução e de dados, o controlador de faltas e o mecanismo de consistência entre os fluxos de instrução e de dados exigido pela instrução FENCE.I.

Dessa forma, o elo entre as duas frentes é a interface, e não a implementação. O contrato de integração estabelece os sinais trocados entre a cache e o pipeline (requisição, endereço, dado de leitura, dado de escrita, indicação de acerto e sinal de stall) e é descrito em termos de largura de bits e de protocolo de troca no \refCap{4_Especificacao}. Em suma, esse contrato permite que os dois subsistemas progridam em paralelo e que cada um seja verificado isoladamente antes da integração.

O código-fonte do projeto está versionado publicamente e é a referência para a descrição apresentada no \refCap{5_Implementacao}~\cite{melo2026riscvchisel}. A versão analisada corresponde ao estado do repositório em 23 de novembro de 2026, quando o pipeline de três estágios e a hierarquia de memória já se encontravam integrados e a bateria de testes automatizados se encontrava em execução estável.

\section{Objetivos}\label{sec:1.3}

O objetivo geral deste trabalho é especificar, implementar e verificar o subsistema de memória cache de um processador RISC-V descrito em Chisel, de modo que o acesso a instruções e a dados deixe de figurar entre os fatores limitantes do desempenho do pipeline.

Para tanto, os objetivos específicos são:

\begin{enumerate}
\def\labelenumi{\arabic{enumi}.}
\item
  Levantar, a partir da literatura, os parâmetros arquiteturais que definem uma memória cache (capacidade, tamanho de linha, grau de associatividade, política de escrita e política de reposição) e justificar as escolhas adotadas para o processador em desenvolvimento.
\item
  Especificar a máquina de estados que governa o atendimento de acertos e faltas, definindo cada estado, as condições de transição e os protocolos de troca com o pipeline e com a memória principal.
\item
  Implementar, em Chisel, as memórias L1 de instrução e de dados como módulos distintos, considerando que a cache de instrução dispensa o bit de sujeira por realizar apenas leituras.
\item
  Implementar o controlador de faltas, responsável por suspender o pipeline durante o atendimento de uma falta, sinalizar o stall ao estágio consumidor e retomar a execução após a conclusão da transferência.
\item
  Implementar o mecanismo de consistência entre os fluxos de instrução e de dados exigido pela instrução FENCE.I, que assegura que instruções gravadas em memória sejam observadas pela cache de instrução na execução subsequente.
\item
  Verificar os módulos por meio de testes automatizados com ScalaTest e ChiselTest sobre o simulador Treadle, comparando os resultados com um modelo de referência.
\item
  Quantificar o efeito da geometria da cache sobre o desempenho, por meio da taxa de acerto e do tráfego de memória em bytes, em diferentes configurações de associatividade, de capacidade, de política de reposição e de tamanho de linha, e determinar o custo em ciclos de cada evento do subsistema.
\end{enumerate}

Os objetivos 1 e 2 correspondem à etapa de especificação, descrita no \refCap{4_Especificacao}. Os objetivos 3 a 5 correspondem aos três componentes que constituem a contribuição delimitada na Seção~\ref{sec:1.2} e são tratados no \refCap{5_Implementacao}. Os objetivos 6 e 7 referem-se à verificação e à avaliação experimental, tratados no \refCap{6_Verificacao}.

\section{Justificativa}\label{sec:1.4}

A justificativa deste trabalho parte de um desequilíbrio persistente na arquitetura de computadores: a frequência de operação dos processadores cresceu continuamente ao longo das últimas décadas, mas a latência de acesso à memória principal não acompanhou esse ritmo. Como consequência, o intervalo entre a emissão de uma requisição de leitura e a disponibilidade do dado passou a dominar o tempo de execução de programas com uso intenso de memória, e o poder de processamento disponível não se converte integralmente em desempenho observável~\cite{wulf1994memorywall}.

Em um processador organizado em pipeline, esse desequilíbrio tem efeito imediato sobre a eficiência. Cada acesso cujo dado não esteja disponível na hierarquia de memória interrompe o avanço das instruções e insere ciclos de espera que não realizam trabalho útil. Portanto, sem um subsistema de cache, o paralelismo que o pipeline oferece é desperdiçado justamente nos trechos de código em que mais haveria ganho a obter. Intervir nesse ponto não é otimização acessória, é a condição para que a organização em estágios produza o ganho que a justifica.

Do ponto de vista do projeto em que este trabalho se insere, a cache ocupa posição de pré-requisito. Uma vez que o objetivo coletivo é a execução do sistema operacional Linux, torna-se necessário um núcleo capaz de sustentar a execução de código compilado, o que envolve desde o carregamento do bootloader até a operação estável do kernel sobre a hierarquia de memória. Nesse contexto, o subsistema de memória não se reduz a acessório de desempenho, ele integra a infraestrutura do projeto, porque, ao manter cópias distintas de instruções e de dados, cria a obrigação de que a consistência entre essas duas visões seja garantida por mecanismo próprio.

Além disso, a escolha do Chisel como linguagem de descrição reforça a pertinência do trabalho em dois aspectos. O primeiro é metodológico, em que, por ser embarcada em Scala, a linguagem permite que os parâmetros da cache (capacidade, tamanho de linha e grau de associatividade) sejam argumentos do gerador, e não constantes repetidas em vários pontos do código. Como resultado, viabiliza a avaliação experimental prevista no \refCap{6_Verificacao} sem que cada configuração exija reescrita. O segundo aspecto é o estado da arte, dado que o ecossistema de projetos RISC-V ainda é majoritariamente descrito em VHDL\glsadd{VHDL} e Verilog, de modo que relatar uma implementação de hierarquia de memória em Chisel, com seu ambiente de verificação automatizada, contribui para a consolidação dessa alternativa na prática de projeto.

Por fim, há uma contribuição de ordem didática e de reprodutibilidade. O subsistema de memória é, entre os componentes de um processador, um dos que mais se beneficiam de documentação detalhada, porque seu comportamento depende de uma máquina de estados cujas transições raramente são visíveis no código final. Logo, ao descrever a especificação, a implementação e a verificação da cache em português, este trabalho disponibiliza material de apoio a disciplinas de arquitetura de computadores e a projetos que venham a reutilizar a hierarquia de memória aqui descrita.
